%=============================================================================
% WAVE 4, ITERATION 11: PATENT PORTFOLIO TRIAGE UPDATE
%
% Agent: Patent Portfolio (W4i11, Opus 4.6)
% Date: 20 March 2026
%
% MANDATE:
%   1. Propagate W4i10 findings through patent triage
%   2. Reassess P4 (65% conversion, $4-12B TAM)
%   3. Reassess P7 (DEW TAM $12.4B-$28B, aircraft-viable)
%   4. Reassess P12 (multi-cavity superradiance 97%)
%   5. Commercial pathway ranking by nearest-to-revenue
%   6. Risk assessment: wide binaries 19sigma, DESI ~3sigma
%
% INPUT:
%   MASTER_FINDINGS_CORPUS.md (W3i8 FINAL)
%   W3i8_Patents_Tech_update.tex (definitive triage)
%   All W4i10_*.tex files (17 documents)
%
% KEY W4i10 INPUTS:
%   P4:  65% conversion at C=1 (omega* ~ 32 GHz), LCOE 19c/kWh, TAM $4-12B
%   P7:  DEW TAM $12.4B (2025)->$28B (2034), aircraft-viable 5.3 m^3
%   P9:  Operationally a gravity gradiometer; mining first revenue $1B
%   P12: Multi-cavity N=3 gives C=1.41, eta=97% (theoretical)
%   P11: SQMS Phase I costed at $3.2M (from $2-4M range)
%   P2:  Dark SRF bound |lambda-1| < 4.2e-10 (3.7x tighter than SQMS)
%   Cosmo: DESI DR2 tension increases to ~2.8-3.2 sigma
%   New: Neutrino mass sum Sigma m_nu = 61.4 meV (zero-parameter)
%   P15: ESS Channel B delayed to ~2029
%=============================================================================

\documentclass[11pt]{article}
\usepackage{amsmath,amssymb,amsthm}
\usepackage{geometry}
\geometry{margin=1in}
\usepackage{booktabs}
\usepackage{enumitem}
\usepackage{hyperref}
\usepackage{xcolor}
\usepackage{tcolorbox}
\usepackage{multirow}
\usepackage{array}
\usepackage{siunitx}
\usepackage{float}
\usepackage{longtable}

\tcbuselibrary{skins,breakable}

\newcommand{\ee}[1]{\times 10^{#1}}
\newcommand{\lambdaone}{|\lambda{-}1|}
\newcommand{\dpsiEM}{\delta\psi_{\rm EM}}
\newcommand{\psigrav}{\psi_{\rm grav}}
\newcommand{\Qpsi}{Q_\psi}
\newcommand{\Meff}{M_{\rm eff}}

\title{\textbf{Wave 4 Iteration 11: Patent Portfolio Triage Update}\\[8pt]
{\large Propagating Iteration 10 Findings Through All 15 Patents}\\[4pt]
{\large Updated Triage Table, Commercial Pathway Ranking, Risk Assessment}\\[4pt]
{\normalsize TOP 5 FINDINGS with Full Justification}}
\author{DFD Research Programme --- Patent Portfolio Agent (W4i11, Opus 4.6)}
\date{20 March 2026}

\begin{document}
\maketitle
\tableofcontents
\newpage

%=============================================================================
\section*{Executive Summary}
%=============================================================================

\begin{tcolorbox}[colback=blue!5!white, colframe=blue!60!black,
  title=\textbf{W4i11: Five Key Findings}]

\textbf{Finding 1 (Triage Change: P7 NICHE $\to$ ALIVE):}
W4i10 established that P7 pulsed energy storage fits the F-35 weapons bay
envelope (5.3~m$^3$, 1550~kg, 46~MJ burst) and addresses a DEW/HPM market
of \$12.4B (2025) growing to \$28B (2034) at 9.5\% CAGR. The three
decisive win conditions (volume-constrained platforms, 4~Hz rep rate at
$>$10~MJ, arbitrary pulse shaping) have no competitor within $10\times$.
Combined TAM across DEW, IFE, and accelerator RF: \$3.2--7.1B.
\textbf{This TAM, the specificity of the competitive advantage, and the
direct military funding pathway ($>\$789$M/yr DoD DE budget) justify
upgrading P7 from NICHE to ALIVE.}

\textbf{Finding 2 (Triage Confirmation: P4 Remains NICHE):}
The frequency-optimised 65\% conversion ($C = 1$ at $\sim$32~GHz) is a
genuine advance, improving LCOE from 27~c/kWh to 19~c/kWh and opening a
\$4--12B TAM. However, P4 remains NICHE because: (a) 19~c/kWh is
competitive only in three restricted market segments (deep submarine,
ultra-deep mining, covert military), not in general power delivery;
(b) the 65\% conversion is a theoretical extrapolation requiring
experimental verification at 32~GHz; (c) all claims remain contingent
on SQMS Phase~I; (d) DFD underground WPT loses to conventional cable
for $<$50~km. P4 is a stronger NICHE than at W3i8, but does not
cross the ALIVE threshold.

\textbf{Finding 3 (Triage Confirmation: P12 Remains THEORETICAL):}
Multi-cavity superradiance ($N = 3$, $\eta = 97\%$) is a theoretical
breakthrough path that would transform P12 economics (LCOE from 27~c/kWh
to $\sim$14~c/kWh). However: (a) the $N^2$ superradiance scaling was
corrected to $NC_1$ in the same document (executive summary was
overclaimed, body corrected it); (b) phase-coherent psi-mode locking
across cavities is undemonstrated; (c) requires SQMS Phase~I as
prerequisite. P12 remains THEORETICAL --- the multi-cavity path is
promising but adds a second layer of unproven physics beyond the
conversion itself.

\textbf{Finding 4 (Nearest-to-Revenue Ranking):}
Updated commercial pathway ranking by time to first revenue:
\begin{enumerate}[nosep]
  \item P9 geological survey service (2033--35): \$500K--2M/campaign.
    Requires SQMS Phase~I + sensor adaptation. Mining TAM \$1B.
  \item P7 DEW/HPM buffer (2034--36): Military procurement pathway.
    Requires SQMS Phase~I + cryo engineering. DEW TAM \$2.5--5.6B.
  \item P4 submarine/mining WPT (2035--38): Requires conversion demo.
    TAM \$4--12B at 19~c/kWh.
  \item P9 submarine navigation (2035--38): \$2--5M/unit.
    TAM \$4--8B. Overlaps P9 geological timeline.
  \item P3 QEC advantage (2036+): Requires quantum hardware maturation.
\end{enumerate}

\textbf{Finding 5 (Risk Assessment: DESI Strengthening Is the Primary
  Threat):}
DESI DR2 increases the CPL tension from 2.4$\sigma$ to $\sim$2.8--3.2$\sigma$.
If DESI DR3 (5-year, $\sim$2027) pushes this above 4$\sigma$, DFD's
dark energy model is in serious jeopardy. This directly affects P5
(cosmological predictions) and indirectly undermines P14 (canonical
framework). Wide binaries at 19$\sigma$ are a pure positive ---
they STRENGTHEN DFD. The dark SRF bound ($\lambdaone < 4.2\ee{-10}$)
tightens the allowed parameter space for SQMS Phase~I, narrowing
the window but not closing it (SQMS Phase~I reach is $7.8\ee{-10}$,
still above the bound). The neutrino mass prediction ($\Sigma m_\nu
= 61.4$~meV) creates a new high-value testable output that could
validate or falsify DFD within 2 years via DESI+Planck.

\end{tcolorbox}

%=============================================================================
%=============================================================================
\part{Updated Patent Triage Table}
\label{part:triage}
%=============================================================================
%=============================================================================

%=============================================================================
\section{Definitive Triage: W4i11 vs.\ W3i8}
\label{sec:triage-table}
%=============================================================================

\begin{table}[H]
\centering
\caption{Patent triage comparison: W3i8 (converged) vs.\ W4i11 (post-W4i10).
  Changes highlighted in \textbf{bold}.}
\label{tab:triage}
\begin{tabular}{@{}clll@{}}
\toprule
Patent & W3i8 Category & W4i11 Category & Change Justification \\
\midrule
P1  & DEAD & DEAD & No change. Signal $10^{15}\times$ below detection. \\
P2  & NICHE & NICHE & Dark SRF tightens bound to $4.2\ee{-10}$. \\
P3  & NICHE & NICHE & No change. 7.7--83$\times$ QEC advantage. \\
P4  & NICHE & NICHE & Strengthened: 65\% conv., LCOE 19~c/kWh, TAM \$4--12B. \\
P5  & ALIVE & ALIVE & Neutrino mass prediction added. DESI tension rising. \\
P6  & NICHE & NICHE & No change. Decihertz gap-filler confirmed. \\
\textbf{P7} & \textbf{NICHE} & \textbf{ALIVE} & \textbf{DEW TAM \$12.4B, aircraft-viable, 3 decisive win conditions.} \\
P8  & DEAD & DEAD & No change. 1.6 Earth masses minimum. \\
P9  & ALIVE & ALIVE & Mining first-revenue pathway defined. Gradiometer risk noted. \\
P10 & NICHE & NICHE & No change. BQP-complete confirmed. \\
P11 & ALIVE & ALIVE & SQMS Phase~I costed at \$3.2M. Dark SRF synergy. \\
P12 & THEORETICAL & THEORETICAL & Multi-cavity path identified but undemonstrated. \\
P13 & NICHE & NICHE & CLONETS-DS protocol observer-ready. \\
P14 & ALIVE & ALIVE & Hidden Axiom 4 assumption noted. T4 ameliorated. \\
P15 & NICHE & NICHE & ESS Channel B delayed to $\sim$2029. \\
\bottomrule
\end{tabular}
\end{table}

\begin{tcolorbox}[colback=green!5!white, colframe=green!60!black,
  title=\textbf{W4i11 Summary: 5 ALIVE, 8 NICHE, 1 THEORETICAL, 2 DEAD}]

\textbf{ALIVE (5)}: P5, P7, P9, P11, P14. \\
\textbf{NICHE (8)}: P2, P3, P4, P6, P10, P13, P15. \\
\textbf{THEORETICAL (1)}: P12. \\
\textbf{DEAD (2)}: P1, P8.

\medskip
\textbf{One change from W3i8}: P7 upgraded NICHE $\to$ ALIVE.

\medskip
\textbf{Near-miss decisions}:
\begin{itemize}[nosep]
  \item P4: Considered for upgrade. Retained as NICHE --- the 19~c/kWh
    LCOE limits viable markets to three niche segments, and the
    65\% conversion requires experimental verification.
  \item P12: Multi-cavity superradiance is promising but both the
    superradiance mechanism and the underlying conversion physics are
    undemonstrated. Two layers of theoretical risk.
  \item P9: Already ALIVE. The honest negative (operationally a gravity
    gradiometer) was considered for downgrade, but the zero-drift,
    covert, all-environment properties remain uniquely DFD-enabled
    at the predicted sensitivity levels.
\end{itemize}
\end{tcolorbox}

%=============================================================================
\section{Detailed Justification: P7 NICHE $\to$ ALIVE}
\label{sec:P7-upgrade}
%=============================================================================

The upgrade of P7 from NICHE to ALIVE is justified by three independent
arguments:

\subsection{Argument 1: Market scale and specificity}

W4i10 established:
\begin{itemize}[nosep]
  \item DEW/HPM global market: \$12.4B (2025), \$28B (2034), CAGR 9.5\%.
  \item HPM is the fastest-growing segment within DEW.
  \item IFE driver market: \$0.5--1.1B (subset of \$40--80B fusion sector).
  \item Accelerator RF: \$0.15--0.4B.
  \item Combined P7 TAM: \$3.2--7.1B (2025), \$6--12B (2034).
\end{itemize}

For comparison, the W3i8 ALIVE patents had the following TAM assessments:
\begin{itemize}[nosep]
  \item P5: No direct TAM (scientific predictions).
  \item P9: \$4--8B (navigation) + \$1B (geological).
  \item P11: No direct TAM (detection gateway experiment).
  \item P14: No direct TAM (theoretical framework).
\end{itemize}

P7's \$3.2--7.1B TAM is comparable to P9's \$5--9B and is the
second-largest commercial opportunity in the portfolio.

\subsection{Argument 2: Decisive competitive advantage}

P7 has three win conditions where no competitor comes within $10\times$:

\begin{enumerate}
\item \textbf{Volume-constrained platforms}: 650~MJ/m$^3$ cavity energy
  density exceeds SMES (best competitor at 57~MJ/m$^3$) by $11\times$.
  The aircraft-viable system (5.3~m$^3$, 1550~kg, 46~MJ) fits the F-35
  weapons bay. No capacitor bank achieves $>$1~MJ in that volume.

\item \textbf{High repetition rate at $>$10~MJ}: 4~Hz vs.\ $<$0.1~Hz for
  capacitors/Marx generators, $\sim$1~Hz for SMES. Factor $4\times$--$400\times$.
  This is decisive for IFE (requires 10~Hz) and rapid-fire DEW.

\item \textbf{Pulse waveform flexibility}: Arbitrary discharge shaping via
  $Q_{\rm ext}$ control with MHz feedback bandwidth. No competing technology
  offers intrinsic shaping at $>$10~MJ without additional pulse-forming
  networks.
\end{enumerate}

\subsection{Argument 3: Direct military funding pathway}

DoD DE budget: \$789.7M (FY2025). This exceeds the total accelerator
RF market. Military procurement is the highest willingness-to-pay
segment and does not require general market LCOE competitiveness.
The path from SQMS Phase~I to a P7 DEW demonstrator is a natural
DoD/DARPA programme.

\subsection{Why this was NICHE at W3i8}

At W3i8, the P7 assessment noted storage time of 1~s, 46~MJ max, and
repositioning as ultra-fast pulsed buffer. The critical missing pieces
that W4i10 supplied:
\begin{itemize}[nosep]
  \item Detailed system design showing aircraft viability (5.3~m$^3$).
  \item Specific market sizing with current (2025--2026) DEW budget data.
  \item Three formally identified win conditions with quantitative
    competitor comparison.
  \item IFE dual-store architecture viability.
  \item FCC-class accelerator \$240--760M lifetime savings.
\end{itemize}

These move P7 from ``viable in restricted domains'' (NICHE) to
``clear path to demonstration and active development warranted'' (ALIVE).

%=============================================================================
\section{Detailed Justification: P4 Remains NICHE}
\label{sec:P4-niche}
%=============================================================================

\subsection{What changed in W4i10}

\begin{center}
\begin{tabular}{lll}
\toprule
Parameter & W3i7/W3i8 & W4i10 \\
\midrule
$\eta_{\rm conv}$ & 50\% & 65\% (at $C = 1$, $\omega^* \approx 32$~GHz) \\
$\eta_{\rm e2e}$ & 18.8\% & 31.8\% \\
LCOE & $\sim$27~c/kWh & $\sim$19~c/kWh \\
TAM & Not quantified & \$4--12B/yr \\
\bottomrule
\end{tabular}
\end{center}

\subsection{Why P4 does not upgrade to ALIVE}

\begin{enumerate}
\item \textbf{Market remains niche}: At 19~c/kWh, P4 is competitive only
  for: (a) deep submarine power ($>$500~km, where HVDC cable $>$\$2.5B),
  (b) ultra-deep mining ($>$3~km), (c) covert military power. It does NOT
  compete for general underground power delivery ($<$50~km) where
  conventional cable is $>$95\% efficient. General viability requires
  LCOE $\sim$10~c/kWh ($\eta_{\rm conv} \sim 80\%$).

\item \textbf{65\% conversion is unverified}: The frequency optimisation
  exploiting $C \propto \omega^{-2}$ is a theoretical extrapolation from
  47~GHz to 32~GHz. The polariton crossing tunability at 32~GHz is assumed
  but not demonstrated. SRF $Q$ at 32~GHz is assumed comparable to 47~GHz
  (reasonable but untested).

\item \textbf{Cascaded conversion remains blocked}: The W4i10 two-frequency
  cascade investigation confirmed that the polariton gap ($\sim 10^{-21}$~Hz)
  precludes branch separation. No route beyond 65\% exists within the DFD
  Lagrangian.

\item \textbf{ALIVE threshold requires ``clear path to demonstration''}:
  P7 has a natural military demonstrator pathway (DEW integration). P4's
  demonstrator would require a full-scale \$270M underground WPT system
  with no intermediate validation step. The step from SQMS Phase~I to
  a 100~MW underground WPT system is too large.
\end{enumerate}

\textbf{Verdict}: P4 is a \emph{strengthened} NICHE. If 65\% conversion is
experimentally confirmed and a military (covert submarine power) customer
emerges, P4 could be reconsidered for ALIVE at a future iteration.

%=============================================================================
\section{Detailed Justification: P12 Remains THEORETICAL}
\label{sec:P12-theoretical}
%=============================================================================

\subsection{The multi-cavity superradiance finding}

W4i10 P12 analysis found that $N = 3$ phase-coherent cavities achieve
$C_3 = 1.41$, giving $\eta = 97\%$ ideal conversion. This would transform
P12 economics:
\begin{itemize}[nosep]
  \item End-to-end: $97\% \times 81\% \times 50\% = 39\%$.
  \item City-scale cost: \$0.3--0.6B (down from \$0.6--1.1B).
  \item LCOE: $\sim$14~c/kWh (approaching general viability).
\end{itemize}

\subsection{Why P12 does not upgrade}

\begin{enumerate}
\item \textbf{Two layers of unproven physics}: P12 requires both
  (a) EM-to-psi conversion (the core bottleneck, gated by SQMS Phase~I)
  AND (b) phase-coherent psi-mode coupling across multiple cavities
  (never attempted, requires psi-mode spatial coherence over $\sim$1~mm
  at 47~GHz).

\item \textbf{Executive summary overclaim corrected in body}: The W4i10
  executive summary stated $C_N = N^2 C_1$ (Dicke superradiance). The
  body of the document corrected this to $C_N = NC_1$ (coupled oscillator).
  The $N = 3$ result holds under the corrected scaling but represents a
  much less dramatic enhancement than initially suggested.

\item \textbf{Receiving-end conversion unchanged}: Even at 97\% transmitting
  conversion, the receiving end remains at 50\% (or 65\% with P4-style
  frequency optimisation). The net efficiency is limited by the weaker end.

\item \textbf{No demonstrator pathway}: There is no intermediate experiment
  between SQMS Phase~I and a multi-cavity superradiant WPT system.
\end{enumerate}

\textbf{Verdict}: P12 remains THEORETICAL. Multi-cavity superradiance is
the most promising conversion breakthrough path and should be prioritised
in the post-SQMS-Phase-I roadmap, but does not change P12's classification
today.

%=============================================================================
%=============================================================================
\part{Commercial Pathway Ranking}
\label{part:commercial}
%=============================================================================
%=============================================================================

%=============================================================================
\section{Nearest-to-Revenue Analysis}
\label{sec:revenue-ranking}
%=============================================================================

All DFD technology patents share a common prerequisite: SQMS Phase~I
detection of $\lambdaone > 0$. Estimated SQMS Phase~I completion:
2027--2028 (\$3.2M, 24 months). All timelines below assume SQMS
Phase~I success.

\begin{table}[H]
\centering
\caption{Patent revenue ranking by estimated time to first revenue.}
\label{tab:revenue-ranking}
\begin{tabular}{@{}clp{4cm}ccc@{}}
\toprule
Rank & Patent & Application & First Revenue & TAM & Confidence \\
\midrule
1 & P9 & Geological survey & 2033--35 & \$1B/yr & Medium-High \\
2 & P7 & DEW/HPM buffer & 2034--36 & \$2.5--5.6B/yr & Medium \\
3 & P4 & Submarine/mining WPT & 2035--38 & \$4--12B/yr & Medium-Low \\
4 & P9 & Submarine navigation & 2035--38 & \$4--8B/yr & Medium \\
5 & P3 & QEC advantage & 2036+ & N/A (enabling) & Low \\
6 & P10 & Psi-field computing & 2038+ & N/A (enabling) & Low \\
7 & P7 & IFE driver & 2038+ & \$0.5--1.1B/yr & Low \\
8 & P7 & Accelerator RF & 2038+ & \$0.15--0.4B/yr & Low \\
\bottomrule
\end{tabular}
\end{table}

%=============================================================================
\section{Rationale for Revenue Ranking}
\label{sec:revenue-rationale}
%=============================================================================

\subsection{Rank 1: P9 Geological Survey (2033--35)}

\textbf{Why first}: Geological survey is the simplest operational mode
--- no real-time navigation required, relaxed precision ($<$1~m), existing
market with established customers. The 23~km depth penetration is
4--5$\times$ beyond airborne gravity gradiometry. The path from SQMS
Phase~I to geological survey requires only sensor adaptation (atom
interferometer to psi-gradient mode), which leverages existing
atom interferometer hardware at TRL~5--6.

\textbf{Honest negative}: W4i10 confirmed that the $\psigrav$-gradient
sensor is operationally identical to a gravity gradiometer at leading
order. The ``psi-field'' framing adds theoretical structure but the
operational device is, practically speaking, a gravity gradiometer.
This means P9's advantage over conventional gravity gradiometry must
come from sensitivity (sub-mm vs.\ $\sim$10~m for airborne FTG) and
depth penetration (23~km vs.\ 5~km), not from any fundamentally
different measurement principle. If conventional quantum gravity
gradiometers achieve comparable sensitivity (currently at TRL~3--4),
P9's competitive moat narrows.

\textbf{Revenue model}: \$500K--2M per geological survey campaign.
Global mineral exploration gravity survey market: $\sim$\$1B/yr.

\subsection{Rank 2: P7 DEW/HPM Buffer (2034--36)}

\textbf{Why second}: Military procurement has the highest willingness-to-pay
and the most direct funding pathway (DoD DE budget \$789.7M/yr). The
5.3~m$^3$/1550~kg system fits existing platforms (F-35). However, the path
from SQMS Phase~I to a flight-qualified DEW system is longer than the
geological survey path: it requires cryogenic engineering for a mobile
platform, integration with existing HPM antennas, and military
qualification testing.

\textbf{Revenue model}: Defence contractor subcontract for pulsed energy
storage subsystem. Per-unit: \$5--20M. Initial procurement: 10--50 units.

\subsection{Rank 3: P4 Submarine/Mining WPT (2035--38)}

\textbf{Why third}: P4 requires conversion efficiency demonstration at
32~GHz (not just detection), plus a \$270M system-level demonstrator.
The market (\$4--12B) is large but the capital requirement for first
installation is prohibitive without government/military anchor customer.
Most likely entry: covert military submarine power delivery, funded by
naval acquisition programmes.

\subsection{Rank 4: P9 Submarine Navigation (2035--38)}

\textbf{Why fourth (not first)}: While submarine navigation has higher
unit value (\$2--5M) and larger TAM (\$4--8B) than geological survey,
the operational requirements are more demanding: real-time, drift-free,
sub-mm precision in a hostile environment. This requires a more mature
sensor than geological survey mode.

%=============================================================================
%=============================================================================
\part{Risk Assessment}
\label{part:risk}
%=============================================================================
%=============================================================================

%=============================================================================
\section{External Risk Factors}
\label{sec:external-risk}
%=============================================================================

\subsection{DESI DR2: Tension Rising to $\sim$3$\sigma$}

\begin{tcolorbox}[colback=red!5!white, colframe=red!60!black,
  title=\textbf{RISK: DESI CPL Tension (W4i10 Cosmo)}]

\textbf{Status}: DESI DR2 (March 2025) reports $w_0 > -1$, $w_a < 0$
preferred at 3.1$\sigma$ (BAO+CMB). DFD predicts $w_0 = -1.183$,
$w_a = +0.183$. The sign of $w_a$ disagrees. Tension rises from
2.4$\sigma$ (DR1) to $\sim$2.8--3.2$\sigma$ (DR2).

\textbf{Impact}: If DESI DR3 ($\sim$2027) pushes above 4$\sigma$, DFD's
dark energy model faces falsification. This directly affects:
\begin{itemize}[nosep]
  \item P5 (cosmological predictions): CPL parameters are part of the
    prediction set.
  \item P14 (canonical framework): If the psi-screen dark energy mechanism
    is falsified, the framework loses one of its cosmological applications
    (but NOT the laboratory/technology predictions).
\end{itemize}

\textbf{Mitigation}: A full DFD Boltzmann code is the critical missing tool.
The CPL parameterisation may be inadequate for DFD's non-standard dark
energy mechanism. A dedicated DFD likelihood analysis against DESI data
could reduce or increase the tension.

\textbf{Patents affected}: P5, P14 (directly). All technology patents
(P2--P4, P6--P7, P9--P13, P15) are \textbf{unaffected} --- they depend on
the EM-psi coupling ($\lambda$), not on the cosmological dark energy model.

\textbf{Severity}: MODERATE-HIGH. Not yet falsifying but trajectory is
concerning.
\end{tcolorbox}

\subsection{Wide Binaries at 19$\sigma$: Pure Positive}

\begin{tcolorbox}[colback=green!5!white, colframe=green!60!black,
  title=\textbf{POSITIVE: Wide Binary Enhancement Confirmed}]

Wide binary observations now show a 19$\sigma$ detection of enhanced
velocities at separations $>$2000~AU, consistent with MOND-like
gravitational modification. DFD predicts +42\% enhancement at 10,000~AU
(M$_\odot$). This is a strong confirmation of DFD's core prediction
(MOND recovery from the interpolating function $\mu(x) = x/(1+x)$).

\textbf{Patents affected}: All patents benefit (framework validation).
Strongest benefit: P14 (canonical framework), P5 (predictions), P9
(gravitational gradient structure used for navigation relies on the
same modified gravity that produces the wide binary signal).

\textbf{Severity}: NONE (purely positive).
\end{tcolorbox}

\subsection{Dark SRF Bound: Tightened to $4.2\ee{-10}$}

\begin{tcolorbox}[colback=yellow!5!white, colframe=yellow!60!black,
  title=\textbf{CAUTION: Dark SRF Narrows SQMS Phase I Window}]

\textbf{Status}: W4i10 P2 analysis reinterprets the Dark SRF null result
at Fermilab as a constraint $\lambdaone < 4.2\ee{-10}$, which is
$3.7\times$ tighter than the SQMS bound of $1.56\ee{-9}$.

\textbf{Impact on SQMS Phase~I}: The Phase~I 4-cavity Q-ratio diagnostic
has reach $\lambdaone \sim 7.8\ee{-10}$. With the dark SRF bound at
$4.2\ee{-10}$, the SQMS experiment must detect a signal in the range
$4.2\ee{-10} < \lambdaone < 7.8\ee{-10}$ (less than a factor of 2).
This is still viable but the window is narrower than assumed at W3i8.

\textbf{Caveat}: The dark SRF bound has a critical uncertainty: the
geometric overlap factor between DFD scalar coupling and dark photon
mixing may weaken the bound by a factor of $\sim$1--10. With a factor-10
weakening: $\lambdaone < 1.3\ee{-9}$, comparable to SQMS.

\textbf{Mitigation}: SQMS Phase~I remains viable. The dark SRF bound
is conservative (likely weakened by geometric factors). Moreover, SQMS
Phase~I's Q-ratio diagnostic is qualitatively different from the dark SRF
power-transfer topology --- a positive SQMS result would be independent
confirmation even if the dark SRF bound holds.

\textbf{Patents affected}: P11 (detection gateway). Indirectly: all
technology patents (gated by P11).

\textbf{Severity}: LOW-MODERATE. Narrows but does not close the window.
\end{tcolorbox}

\subsection{ESS Channel B: Delayed to $\sim$2029}

The ESS moderator setback delays Channel B to $\sim$2029 (from 2028--30).
Physics case unchanged. Impact: P15 (passive detection programme) and P11
(detection roadmap) timelines extended by $\sim$1 year.

\textbf{Severity}: LOW. Timeline shift only.

\subsection{Neutrino Mass Prediction: New High-Value Output}

\begin{tcolorbox}[colback=green!5!white, colframe=green!60!black,
  title=\textbf{OPPORTUNITY: $\Sigma m_\nu = 61.4$~meV --- Testable by 2027}]

W4i10 identified a zero-parameter prediction from the DFD microsector:
$\Sigma m_\nu = 61.4$~meV, with $(m_1, m_2, m_3) = (2.34, 8.96, 50.12)$~meV
(normal ordering required).

\textbf{Tests}:
\begin{itemize}[nosep]
  \item DESI+Planck: $\Sigma m_\nu$ constraint, 2026--2027.
  \item JUNO: Mass ordering determination, 2027.
  \item Euclid+CMB-S4: $\sigma(\Sigma m_\nu) \sim 15$--20~meV, 2027--2029.
\end{itemize}

\textbf{Falsification}: $\Sigma m_\nu < 45$~meV or $> 80$~meV, or
inverted ordering confirmed.

\textbf{Impact}: This is the most testable new DFD prediction. A
confirmation would be transformative for the entire programme. A
falsification would challenge the microsector derivation but would NOT
affect the 4-axiom framework or technology patents (the neutrino mass
prediction uses the microsector integers, which are independent of the
EM-psi coupling $\lambda$).

\textbf{Patents affected}: P14 (framework validation). Technology patents
unaffected.
\end{tcolorbox}

%=============================================================================
\section{Risk Matrix: Which Patents Are Affected by External Factors?}
\label{sec:risk-matrix}
%=============================================================================

\begin{table}[H]
\centering
\caption{Risk matrix: external factors vs.\ patent impact.}
\label{tab:risk-matrix}
\begin{tabular}{@{}lcccccc@{}}
\toprule
& DESI & Wide & Dark SRF & ESS & $\nu$-mass & T4 Cold \\
& $\sim$3$\sigma$ & 19$\sigma$ & $4.2\ee{-10}$ & delay & 61.4 meV & Spot \\
\midrule
P1 (DEAD) & --- & --- & --- & --- & --- & --- \\
P2 & --- & --- & \textbf{Mod} & --- & --- & --- \\
P3 & --- & --- & Low & --- & --- & --- \\
P4 & --- & --- & Low & --- & --- & --- \\
P5 & \textbf{High} & \textcolor{green!50!black}{++} & --- & --- & \textcolor{green!50!black}{++} & \textbf{Mod} \\
P6 & --- & --- & Low & --- & --- & --- \\
P7 & --- & --- & Low & --- & --- & --- \\
P8 (DEAD) & --- & --- & --- & --- & --- & --- \\
P9 & --- & \textcolor{green!50!black}{+} & Low & --- & --- & --- \\
P10 & --- & --- & Low & --- & --- & --- \\
P11 & --- & --- & \textbf{Mod} & Low & --- & --- \\
P12 & --- & --- & Low & --- & --- & --- \\
P13 & --- & --- & --- & --- & --- & --- \\
P14 & \textbf{Mod} & \textcolor{green!50!black}{++} & --- & --- & \textcolor{green!50!black}{++} & \textbf{Mod} \\
P15 & --- & --- & --- & \textbf{Mod} & --- & --- \\
\bottomrule
\end{tabular}
\end{table}

\textbf{Key}: High = directly threatens patent viability. Mod = affects
scope or timeline. Low = indirect through SQMS gateway. $+/++$ = positive
impact. --- = no impact.

\textbf{Interpretation}: The DESI tension is the single most important
external risk factor. It threatens P5 and P14 (cosmological predictions
and framework), but does NOT affect the technology patents (P2--P4,
P6--P7, P9--P13, P15), which depend on $\lambda$ coupling, not on dark
energy. The technology patents are insulated from cosmological risks.

%=============================================================================
%=============================================================================
\part{Updated Detection Roadmap}
\label{part:roadmap}
%=============================================================================
%=============================================================================

\begin{table}[H]
\centering
\caption{Detection roadmap: W4i11 update (changes from W3i8 in \textbf{bold}).}
\label{tab:roadmap}
\begin{tabular}{@{}clccl@{}}
\toprule
Rank & Experiment & Cost & Reach $\lambdaone$ & Timeline \\
\midrule
\textbf{0} & \textbf{Dark SRF reanalysis} & \textbf{\$0} & \textbf{$4.2\ee{-10}$} & \textbf{Now} \\
1 & SQMS Phase~I (4-cavity) & \textbf{\$3.2M} & $7.8\ee{-10}$ & 2027--28 \\
2 & LIGO archival reanalysis & \$0.1--0.5M & $\sim 3\ee{-14}$ & Now \\
3 & Channel B (ESS) & \$2--5M & $\sim 8\ee{-13}$ & \textbf{$\sim$2029} \\
\bottomrule
\end{tabular}
\end{table}

\textbf{Changes from W3i8}:
\begin{enumerate}[nosep]
  \item \textbf{NEW Rank 0}: Dark SRF reanalysis at zero cost. This is the
    cheapest possible improvement to the $\lambda$ bound.
  \item SQMS Phase~I cost refined from \$2--4M to \$3.2M (W4i10 full
    costing).
  \item ESS Channel B delayed from 2028--30 to $\sim$2029 (moderator setback).
\end{enumerate}

%=============================================================================
%=============================================================================
\part{Summary and Recommendations}
\label{part:summary}
%=============================================================================
%=============================================================================

\section{Changes to MASTER\_FINDINGS\_CORPUS}

\begin{table}[H]
\centering
\caption{W4i11 updates to the Master Findings Corpus.}
\label{tab:corpus-updates}
\begin{tabular}{@{}lp{5cm}p{5cm}@{}}
\toprule
Item & Previous (W3i8/W4i10) & Updated (W4i11) \\
\midrule
Patent triage & 4 ALIVE, 8 NICHE, 1 THEO, 2 DEAD &
  \textbf{5 ALIVE}, 8 NICHE, 1 THEO, 2 DEAD \\
P7 classification & NICHE & \textbf{ALIVE} \\
SQMS Phase~I cost & \$2--4M & \$3.2M (confirmed) \\
$\lambda$ bound & $1.56\ee{-9}$ (SQMS) &
  $4.2\ee{-10}$ (Dark SRF, pending geometric factor) \\
DESI tension & 2.4$\sigma$ (DR1) & $\sim$2.8--3.2$\sigma$ (DR2) \\
Neutrino mass & Not highlighted & $\Sigma m_\nu = 61.4$~meV (prediction 15) \\
ESS Channel B & 2028--30 & $\sim$2029 \\
\bottomrule
\end{tabular}
\end{table}

\section{Priority Actions}

\begin{enumerate}
\item \textbf{Highest priority}: Secure SQMS Phase~I funding (\$3.2M).
  This remains the single gateway experiment for all technology patents.

\item \textbf{Dark SRF geometric factor calculation}: Determine whether the
  $4.2\ee{-10}$ bound holds or weakens under the correct DFD overlap integral.
  Cost: \$0 (theoretical calculation). Impact: narrows or widens the
  SQMS Phase~I discovery window.

\item \textbf{DFD Boltzmann code}: Build the full Boltzmann code for
  DESI/Euclid likelihood analysis. This is the only way to determine
  whether the DESI tension is a real falsification or a parameterisation
  artefact. Cost: 6--12 months of dedicated effort.

\item \textbf{P7 DEW concept study}: Prepare a DARPA/ONR briefing on
  the P7 pulsed energy buffer for HPM applications. Target: \$5--10M
  Phase~II DEW demonstrator funding, contingent on SQMS Phase~I.

\item \textbf{P9 geological MVP}: Define the minimum viable atom
  interferometer configuration for psi-gradient geological survey.
  Prepare white paper for mining industry partners.
\end{enumerate}

\section{Convergence Assessment}

W4i11 produces one triage change (P7: NICHE $\to$ ALIVE) from W3i8.
All other classifications are confirmed. The portfolio is now at
5 ALIVE / 8 NICHE / 1 THEORETICAL / 2 DEAD.

The single triage change is driven by new quantitative market and
system-design data from W4i10, not by any change in the underlying
physics. No W4i10 finding contradicts any W3i8 conclusion. The
correction rate continues to decline (W4i10 produced 0 corrections
to inherited facts for P7; the only corrections were the IFE
dual-store refinement and the P12 superradiance $N^2 \to N$
correction, both internal to the iteration).

The portfolio assessment is converging. Further iteration is warranted
only if external data changes (DESI DR3, SQMS Phase~I results, dark
SRF geometric factor calculation).

\end{document}
